% Lösungen Schnittmengen (zusammenbau v0.4, Heft)
\einheitenkopf{Schnittmengen}

\erg{57}{a) einsetzen – die Geradenkoordinaten in die Koordinatengleichung. \quad b) Einsetzen: $(-1 + t) + 2 + (1 + t) = 6$, $2 + 2t = 6$, also $t = 2$; Schnittpunkt $S(1|2|3)$}
\erg{58}{a) Richtungsvektor $\overrightarrow{PM} = (3|1|-3)$; $\vec{x} = (0|0|3) + t \cdot (3|1|-3)$ in $L$: $3t + 2t + 3 - 3t = 4$, also $t = 0{,}5$ und $Q(1{,}5|0{,}5|1{,}5)$ \quad b) Senkrechte Gerade $\vec{x} = (4|6|2) + t \cdot (0|0|1)$ in $F$: $8 + 18 + 5(2 + t) = 51$, also $t = 3$; Länge $3 \cdot 50 + 1{,}5 = 151{,}5$ m \quad c) Lichtrichtung $\overrightarrow{AA'} = (-2|1|-2)$; Lichtgerade $\vec{x} = (-2|4|4) + t \cdot (-2|1|-2)$; $4 - 2t = 0$ liefert $t = 2$ und $C'(-6|6|0)$ \quad d) Lichtgerade $\vec{x} = (0|5|9) + t \cdot (2|-1|-2)$; $9 - 2t = 6$ liefert $t = 1{,}5$ und den Punkt $(3|3{,}5|6)$ \quad e) Lichtrichtung $\overrightarrow{PP'} = (1|-1|-2)$; Lichtgerade $\vec{x} = (-1|1|5) + t \cdot (1|-1|-2)$; $5 - 2t = 0$ liefert $t = 2{,}5$ und $S'(1{,}5|-1{,}5|0)$ \quad f) Bahn $\vec{x} = (6|3|0) + r \cdot (-2|1|0)$, der Schritt $r = 1$ dauert 4 s; in $E$ eingesetzt ergibt sich $r = 1{,}5$; Zeit $r \cdot 4 = 6$ s \quad g) Bahn $\vec{x} = (4|6|0) + r \cdot (3|-2|0)$, der Schritt $r = 1$ dauert 2 min; in $E$ eingesetzt ergibt sich $r = 1{,}5$; Zeit $r \cdot 2 = 3$ min \quad h) Bahn $\vec{x} = (3|3|5) + r \cdot (1|2|1)$, der Schritt $r = 1$ dauert 5 s; in $E$ eingesetzt ergibt sich $r = 3$; Zeit $r \cdot 5 = 15$ s}
\erg{59}{a) drei Gleichungen mit zwei Unbekannten – eine je Koordinate, unbekannt sind nur $r$ und $s$. \quad b) Gleichsetzen: $2 + r = 3$, $2 = s$, $3 + 2r = 3 + s$; aus zwei Gleichungen $r = 1$, $s = 2$; die dritte stimmt als Probe; Schnittpunkt $S(3|2|5)$}
\erg{60}{a) Gleichsetzen mit $\vec{x} = (8|0|1) + \mu \cdot (0|6|0)$: $2 + 4\lambda = 8$ liefert $\lambda = 1{,}5$; $4{,}5 - \lambda = 6\mu$ liefert $\mu = 0{,}5$; Probe mit der dritten Koordinate: $0{,}25 + 0{,}5 \cdot 1{,}5 = 1$ stimmt; $\mu$ liegt in $[0; 1]$, also $Q(8|3|1)$ \quad b) $h: \vec{x} = (3|0|1) + r \cdot (1|1|b - 1)$; gleichsetzen mit verschiedenen Buchstaben: aus zwei Koordinaten $r = 4$ und $s = 3$; die dritte liefert $b = 2{,}5$ \quad c) $h: \vec{x} = (2|1|0) + r \cdot (1|1|b)$; gleichsetzen mit verschiedenen Buchstaben: aus zwei Koordinaten $r = -3$ und $s = -1$; die dritte liefert $b = -2$ \quad d) $h: \vec{x} = (1|3|2) + r \cdot (1|0|b - 2)$; gleichsetzen mit verschiedenen Buchstaben: aus zwei Koordinaten $r = 2$ und $s = 1$; die dritte liefert $b = 2{,}5$ \quad e) Kante $GH$: $\vec{x} = (0|4|4) + w \cdot (4|0|0)$ mit $0 \le w \le 1$. Gleichsetzen: die zweite Koordinate liefert $u = 1$, die dritte $r^2 = 4$, also $r = 2$ oder $r = -2$; die erste liefert für $r = 2$ den Wert $w = 0{,}25$, für $r = -2$ den Wert $w = 1{,}25$ – außerhalb der Kante, verworfen. Der Schnittpunkt $Q(1|4|4)$ teilt $HG$ im Verhältnis $HQ : QG = 0{,}25 : 0{,}75 = 1 : 3$. \quad f) Kante $GH$: $\vec{x} = (0|6|6) + w \cdot (6|0|0)$ mit $0 \le w \le 1$. Gleichsetzen: die zweite Koordinate liefert $u = 2$, die dritte $r^2 - 3 = 6$, also $r = -3$ oder $r = 3$; die erste liefert für $r = -3$ den Wert $w = 0{,}3333$, für $r = 3$ den Wert $w = 1{,}3333$ – außerhalb der Kante, verworfen. Der Schnittpunkt $Q(2|6|6)$ teilt $HG$ im Verhältnis $HQ : QG = 0{,}3333 : 0{,}6667 = 1 : 2$. \quad g) Kante $GH$: $\vec{x} = (0|2|2) + w \cdot (2|0|0)$ mit $0 \le w \le 1$. Gleichsetzen: die zweite Koordinate liefert $u = 1$, die dritte $r^2 - 2 = 2$, also $r = 2$ oder $r = -2$; die erste liefert für $r = 2$ den Wert $w = 0{,}5$, für $r = -2$ den Wert $w = 1{,}5$ – außerhalb der Kante, verworfen. Der Schnittpunkt $Q(1|2|2)$ teilt $HG$ im Verhältnis $HQ : QG = 0{,}5 : 0{,}5 = 1 : 1$. \quad h) Kante $GH$: $\vec{x} = (0|8|8) + w \cdot (8|0|0)$ mit $0 \le w \le 1$. Gleichsetzen: die zweite Koordinate liefert $u = 2$, die dritte $2r^2 = 8$, also $r = -2$ oder $r = 2$; die erste liefert für $r = -2$ den Wert $w = 0{,}625$, für $r = 2$ den Wert $w = 1{,}625$ – außerhalb der Kante, verworfen. Der Schnittpunkt $Q(5|8|8)$ teilt $HG$ im Verhältnis $HQ : QG = 0{,}625 : 0{,}375 = 5 : 3$.}
\erg{61}{a) Lichtgerade $\vec{x} = (0|0|7) + r \cdot (4|-1|-2)$ und Kantengerade $\vec{x} = (6|-3|4) + s \cdot (0|6|0)$ gleichsetzen: $r = 1{,}5$, $s = 0{,}25$, die dritte Gleichung stimmt; Schattenpunkt $(6|-1{,}5|4)$ – er liegt auf der Kante, weil $s$ zwischen 0 und 1 liegt. \quad b) Lichtgerade $\vec{x} = (1|1|7) + r \cdot (2|1|-3)$ und Kantengerade $\vec{x} = (3|-1|4) + s \cdot (0|6|0)$ gleichsetzen: $r = 1$, $s = 0{,}5$, die dritte Gleichung stimmt; Schattenpunkt $(3|2|4)$ – er liegt auf der Kante, weil $s$ zwischen 0 und 1 liegt. \quad c) Lichtgerade $\vec{x} = (0|0|4{,}5) + r \cdot (-2|2|-1)$ und Kantengerade $\vec{x} = (-3|1|4) + s \cdot (4|0|0)$ gleichsetzen: $r = 0{,}5$, $s = 0{,}5$, die dritte Gleichung stimmt; Schattenpunkt $(-1|1|4)$ – er liegt auf der Kante, weil $s$ zwischen 0 und 1 liegt.}
\erg{62}{a) $x$ und $z$ – auf der $y$-Achse sind die beiden anderen Koordinaten null. \quad b) $S_x(3|0|0)$, $S_y(0|2|0)$, $S_z(0|0|1)$ – je zwei Koordinaten null gesetzt; die Achsenabschnitte sind $\frac{6}{2}$, $\frac{6}{3}$ und $\frac{6}{6}$, nicht die Koeffizienten.}
\erg{63}{a) $S_x(4|0|0)$, $S_y(0|6|0)$, $S_z(0|0|8)$ – je zwei Koordinaten null gesetzt; die Achsenabschnitte sind $\frac{24}{6}$, $\frac{24}{4}$ und $\frac{24}{3}$, nicht die Koeffizienten. \quad b) $F(0|0|4)$ aus $x = y = 0$; die Spurgerade in der $xz$-Ebene verbindet $B(5|0|0)$ und $F$, die in der $yz$-Ebene verläuft wegen $z = 4$ durch $F$ und $E(0|5|4)$. \quad c) $F(0|0|3)$ aus $x = y = 0$; die Spurgerade in der $xz$-Ebene verbindet $B(7|0|0)$ und $F$, die in der $yz$-Ebene verläuft wegen $z = 3$ durch $F$ und $E(0|7|3)$. \quad d) Die Seite der Schnittfigur in der Seitenfläche $x = 7$ (von $A$ nach oben) über die Deckfläche hinaus verlängern, bis sie $h$ trifft; dieser Schnittpunkt ist $P$. Seine Höhe abzählen: $z = 4$ – $P$ liegt nicht in der Schnittfigur selbst. \quad e) Die Seite der Schnittfigur in der Seitenfläche $x = 5$ (von $A$ nach oben) über die Deckfläche hinaus verlängern, bis sie $h$ trifft; dieser Schnittpunkt ist $P$. Seine Höhe abzählen: $z = 4$ – $P$ liegt nicht in der Schnittfigur selbst.}

63b)

\begin{ksys3}[xyz,x1min=0,x1max=6,x2min=0,x2max=6,x3min=0,x3max=5] \rpyramide{0,0,0}{5}{5}{0,5,4} \rpunkt*{0}{0}{4}{F} \rgerade[0:1]{5,0,0}{-5,0,4}{} \rgerade[0:1]{0,0,4}{0,5,0}{} \end{ksys3}


63c)

\begin{ksys3}[xyz,x1min=0,x1max=8,x2min=0,x2max=8,x3min=0,x3max=4] \rpyramide{0,0,0}{7}{7}{0,7,3} \rpunkt*{0}{0}{3}{F} \rgerade[0:1]{7,0,0}{-7,0,3}{} \rgerade[0:1]{0,0,3}{0,7,0}{} \end{ksys3}

\erg{64}{a) $E_1$ schneidet das Quadrat in der Strecke mit $x = 4$ – zwei gleich große Rechtecke, ja. $E_2$ enthält $B$ und $D$ und schneidet entlang der Diagonale $BD$ – zwei kongruente Dreiecke, ja. $E_3$ schneidet bei $y = 2$ – Rechtecke mit den Breiten 2 und 6, nein. \quad b) $E_1$ trifft das Quadrat nicht, denn dort ist $z = 0$ – nein. $E_2$ enthält $A$ und $C$ und schneidet entlang der Diagonale $AC$ – zwei kongruente Dreiecke, ja. $E_3$ schneidet bei $x = 2$ – Rechtecke mit den Breiten 2 und 3, nein. \quad c) $E_1$ schneidet bei $y = 4$, genau in der Mitte zwischen 1 und 7 – ja. $E_2$ enthält $B$ und $D$ und schneidet entlang der Diagonale $BD$ – ja. $E_3$ schneidet bei $x = 5$ – Rechtecke mit den Breiten 4 und 2, nein. \quad d) Ein Normalenvektor von $E_1$ ist $(1|1|1)$ (senkrecht zu beiden Spannvektoren), einer von $E_2$ ist $(2|2|1)$; sie sind keine Vielfachen voneinander, also sind die Ebenen nicht parallel. Parameterform von $E_1$ in $E_2$ einsetzen: $2(5 - r - s) + 2r + s = 10$, $10 - s = 10$, also $s = 0$; Schnittgerade $\vec{x} = (5|0|0) + r \cdot (-1|1|0)$ \quad e) Ein Normalenvektor von $E_1$ ist $(3|2|6)$ (senkrecht zu beiden Spannvektoren), einer von $E_2$ ist $(2|1|4)$; sie sind keine Vielfachen voneinander, also sind die Ebenen nicht parallel. Parameterform von $E_1$ in $E_2$ einsetzen: $2 \cdot 2r + (3 - 3r - 3s) + 4s = 4$, $3 + r + s = 4$, also $s = 1 - r$; Schnittgerade $\vec{x} = (0|0|1) + r \cdot (2|0|-1)$}
